\clearpage
\section{Verification: what each check actually establishes}\label{sec:verification}
\subsection{A chain of evidence, not a count of passing tests}
For a final claim to be useful, the implementation must compute the intended
quantity, use the intended inputs, preserve the comparison, and produce records
that can be checked independently. A large software test suite supports that
work, but its size is not a measure of scientific completion.

\begin{longtable}{@{}P{0.24\linewidth}P{0.35\linewidth}P{0.35\linewidth}@{}}
\toprule Question & Evidence that addresses it & What it does not establish\\\midrule\endhead
Are the equations implemented correctly? & Independent hand calculations of
forward values and the declared backward rules, including edge cases & That the
method works on real movements.\\\addlinespace
Does the learning path work? & A deliberately learnable synthetic problem with
healthy gradients and correct self-excluded retrieval & Good generalization or
real-data superiority.\\\addlinespace
Are the intended examples used? & Independently regenerated manifests, source
identity, overlap/exclusion checks & Independence of every observation.\\\addlinespace
Does the adapter preserve intended encoder behavior? & Local and upstream
encoder outputs on the same prepared input arrays & Byte-identical upstream
preprocessing where the local confidence rule intentionally differs.\\\addlinespace
Does an execution repeat? & Fresh-process initialization, geometry, and full-cache
repeats under the bound runtime & Correctness merely from repeating the same
error, or equality across arbitrary hardware.\\\addlinespace
Was the intervention applied? & Fixed-panel geometry and complete-query damage
summaries & Better robustness of any learned recipe.\\\addlinespace
Can results be reconstructed? & Independent projection, ranking, metric, and
report replay from bound records & Permission to reinterpret the locked claim.\\\addlinespace
Can the whole study finish? & Measured end-to-end costs, retained storage, job
windows, queue allowance, and researcher availability & Feasibility inferred
from fast training alone.\\\bottomrule
\end{longtable}

\subsection{Independent checks must be capable of disagreeing}
Calling the same production function twice can establish repeatability, but
both calls share its possible bug. A mathematical reference should derive a
small result independently, with enough simplicity to inspect manually.
For the contextual loss, both the forward formula and prescribed STE backward
rule need such evidence. A second implementation copied line for line would
not provide the same independence.

Likewise, a training process writing ``complete'' is not independent proof that
the checkpoint, inputs, and scores agree. A separate validation process checks
the retained outputs and their dependencies. A clean process also prevents an
unrecorded in-memory state from supplying missing information.

Real pilots check finite loss, gradients, and parameters, nonzero projections,
unit-vector normalization, adequate embedding variation, complete update counts,
paired initialization/batches, and measured cost. The collapse check asks whether
every example has become essentially the same vector; it does not certify that
the remaining variation is useful. There is no real-data accuracy floor that
forces the experiment to produce a favorable result.

\subsection{The reproducibility problem encountered here}
Floating-point operations round their results. Reordering sums, choosing a
different vectorized kernel, or changing the numerical runtime can change bits
even when the source expression is unchanged. A fixed random seed governs the
random-number sequence; it does not fix every arithmetic path or transformation
of those random numbers.

This project encountered both starting-weight reconstruction mismatches and
small geometry differences across execution environments. In the geometry case,
clean arrays and sample identities matched, while some derived torso scales
and jitter arrays differed. Exact repeatability within each environment did not
establish equality between them. Several diagnostics narrowed possibilities,
but the complete historical numerical cause was not established.

The authorized response was to select and record one prospective runtime,
require exact fresh-process repeatability and pairing within it, and preserve
differences from older outputs. That is an explicit change in the acceptance
requirement. It is not a claim that the old and new numbers were equal after all.

A later cluster kernel update failed the frozen-runtime startup check before
scientific training began. Such a startup failure is evidence that the runtime
binding changed, not evidence that the model produced incorrect numbers on the
new kernel. The new profile was authorized for verification, but compatibility
had to be measured rather than assumed.

\subsection{Immutable evidence and the scope of a correction}
An artifact is retained with its source, input, configuration, and runtime
bindings. A corrected release writes new artifacts with explicit predecessor
links instead of overwriting the failed attempt. This preserves the distinction
between what originally happened and what a later correction established.

The relevant software-engineering idea is dependency-aware invalidation. A
changed scientific input or implementation invalidates evidence depending on
it. It need not invalidate an unrelated unchanged manifest, but old numerical
outputs cannot be relabeled as new-runtime results without the required check.
The protocol determines which dependencies and acceptance evidence matter.

Changes are recorded with a reason and their effect on scientific settings.
The runtime amendments arose from observed reproducibility and execution
problems, with earlier development outputs already in existence. They were not
claimed to predate all development scores, and final novel outcomes remained
unavailable. The algorithms, action IDs, corruption definitions, training
duration, and claim rule stayed fixed.

\subsection{Verification itself has a cost}
The study discovered that nested validation repeatedly recomputed the same
expensive numerical checks. Under the measured old execution path, one repeated
component alone forecast over 122 hours against the old execution plan's shared
48-hour maximum-job check. That was a forecast from measured component time and call counts, not an
observed 122-hour run.

The authorized optimization reuses a successful numerical verification only
inside one guarded command, after rehashing its complete bound contents and
matching the validator settings and identities. A changed dependency prevents
reuse. Independent commands begin without inherited proof. Authorization,
dataset-opening state, resources, and other dynamic checks remain fresh.

This differs from simply skipping expensive checks: the question is whether the
\emph{identical deterministic proposition} has already been established in the
current command. The design still requires independent recertification and new
measurements of total cost.

Storage feasibility also includes retained attempts, checkpoints, feature
caches, projections, logs, temporary peaks, and a reserve. The declared rule is
\[
 \min(\text{filesystem free},\text{quota remaining})
 \ge 200\cdot2^{30}+1.25\cdot\text{remaining peak bytes}.
\]
Free disk space outside the account's quota is not usable capacity. Nor is
deleting failed scientific evidence an acceptable way to make the forecast fit.
Time forecasts similarly need the complete workflow, queue/setup allowances,
and researcher effort, not just the head-training loop.

\source{\src{docs/decisions/0006-v3-execution-completion-amendment.md}{Execution-completion amendment};
\src{docs/decisions/0007-safe-validation-reuse-and-stage-capacity.md}{Validation reuse and capacity};
\src{docs/decisions/0008-runtime-baseline.md}{Prospective runtime decision};
\src{docs/decisions/0009-updated-kernel-baseline.md}{Updated-kernel decision}.}

\section{Decisions and progress through the evidence cutoff}\label{sec:history}
\subsection{The decisions that explain the current experiment}
\begin{longtable}{@{}P{0.19\linewidth}P{0.37\linewidth}P{0.38\linewidth}@{}}
\toprule Decision & Reason & Consequence for the claim or process\\\midrule\endhead
Choose one-shot motion retrieval & Make movement comparison testable on an
existing dataset with unfamiliar action labels & Evaluate matching to one clean
example per novel action; do not claim a rowing-quality system.\\\addlinespace
Adopt a verified pose input contract & Available HRNet-derived annotations and
documented missing captures differed from initial assumptions & Bind the actual
source inventory and regenerate exact manifests.\\\addlinespace
Freeze MotionBERT and narrow to three recipes (29 Sep) & Bound the December
question and resource requirements & Compare head-training recipes; preserve
the broader fine-tuned v2 benchmark as unfinished historical scope.\\\addlinespace
Use primary and official query sets & Synchronized views can duplicate the
gallery performance & Primary excludes all such views; official-compatible
results remain descriptive.\\\addlinespace
Fix duration, pairing, grid, and claim rule & Limit unequal opportunity and
outcome-driven choices & Twenty epochs, three paired seeds, 27 selection runs,
nine final heads, one confirmatory effect.\\\addlinespace
Omit the extra regularizer control & Keep the December study bounded before
new selection outcomes & A result supports the full recipe, not the isolated
neighborhood mechanism.\\\addlinespace
Complete and rehearse the full path (9 Oct) & Successful pilots did not establish
that authorization through report generation worked & Successor releases and
fresh affected evidence before selection.\\\addlinespace
Permit guarded verification reuse (12 Oct) & Measured repeated validation exceeded
an execution limit & Preserve checks while avoiding duplicate numerical work
within an identical guarded command.\\\addlinespace
Adopt one prospective runtime (19 Oct) & Historical cross-environment equality
was not established & Keep exact within-runtime repeatability; record historical
differences honestly.\\\addlinespace
Revalidate after kernel update (25 Oct) & The observed runtime no longer matched
the accepted profile & Authorize a new profile, preserve source/settings, require
fresh affected evidence.\\\bottomrule
\end{longtable}

\subsection{What has actually been completed}
The early work established the literature basis, prospectus, dataset/access
arrangements, verified input inventory, deterministic split manifests, and
adapter implementation. September extraction established encoder agreement on
the same prepared arrays and two complete, matching 95,001-row feature passes.
Those remain historical evidence for their recorded release.

On October 8, Week 4 closed with fresh v3 caches and all three complete
twenty-epoch engineering pilots. The two full extraction passes matched exactly,
and independent checks accepted the outputs. Each pilot had 47,520 updates with
paired initial weights and batches. These are substantial execution milestones,
but the pilot scores are ineligible for selecting a winner.

October successor work completed execution/report rehearsal and the historical
full-query damage audit. The audit checked nine corruption conditions across
18,929 development queries each. Later work accepted a prospective CPU baseline,
its repeatable initialization and geometry, a synthetic GPU prerequisite, and
five caches under the then-frozen profile. The subsequent kernel change means
those records do not automatically satisfy the new-profile requirements.

The project also has a public tracker, research journal, explanatory pose
visualizations, and report/poster drafts with results explicitly pending.
Those help communicate and retain the work; they are not additional scientific
evidence of recipe superiority.

\subsection{What remains incomplete}
At the October 28 checkpoint, Week 5 remains open and cluster jobs are held
while adequate usable storage is pending. The local launcher correction and
conditional remaining-work schedule are prepared. No new-profile acceptance
or final comparison is inferred from that preparation.

The next scientific prerequisites are accepted new-runtime repeatability,
affected caches and paired pilots, development corruption/retrieval replay,
and a measured full-workload capacity decision. Only after admission can the
27 selection runs proceed, followed by nine final heads, final locks, and the
novel evaluation campaign. Selection, final training, and novel opening remain
closed in the latest record.

The original October 18 admission deadline was missed. November 29 final locks,
December 6 computation, and December 7--18 analysis/delivery remain recorded
targets; the guide does not claim the remaining schedule has been demonstrated
feasible. Live capacity was not remeasured while preparing this guide.

\begin{concept}{The accurate status sentence}
The study design and software have passed substantial engineering checks,
including historical full pilots and damage audits. The latest runtime still
needs accepted regeneration, and storage is blocking the next cluster work.
There is no final evidence yet that contextual training is better, equivalent,
or worse on the declared novel-action comparison.
\end{concept}

\source{\src{docs/protocol/week-3-scc-execution.md}{September execution evidence};
\src{docs/protocol/week-4-execution.md}{Week 4 closure};
\src{docs/protocol/week-5-execution.md}{October accepted history and kernel/storage blocker};
\prepsrc{docs/protocol/week-5-execution.md}{October 28 checkpoint};
\prepsrc{docs/protocol/week-6-handoff.md}{Conditional remaining-work calendar}.}
